\documentclass[10pt,a4paper]{amsart}
\usepackage[margin=19mm]{geometry}
\usepackage{amsmath,amssymb,mathtools}
\usepackage[scr=rsfs,cal=euler]{mathalfa}
\usepackage{xfrac}
\usepackage[unicode,hidelinks,bookmarks=false]{hyperref}
\newcommand{\ie}{i.e.\ }
\newcommand{\cf}{cf.\ }
\newcommand{\resp}{respectively\ }
\newcommand{\Subsection}{\subsection}
\providecommand{\nobreakdash}{\nobreak\hbox{-}}
\setlength{\emergencystretch}{2em}
\multlinegap0pt
\usepackage{amssymb}
\newtheorem{thm}{Theorem}
\newtheorem{cor}{Corollary}
\newtheorem{pf}{Pf}
\newcommand{\IC}{{\mathbb C}}
\title{The Landau-Bloch type theorems for certain class of holomorphic and pluriharmonic mappings in $\mathbb{C}^n$}
\author{Vasudeva Rao Allu, Rohit Kumar}
\date{Source: arXiv:2308.15913v4 (CC BY 4.0)}
\begin{document}
\maketitle

\noindent\emph{Source and licence.} The Landau-Bloch type theorems for certain class of holomorphic and pluriharmonic mappings in $\mathbb{C}^n$. Version arXiv:2308.15913v4 (CC BY 4.0), distributed by arXiv under the Creative Commons Attribution 4.0 International licence, http://creativecommons.org/licenses/by/4.0/, as recorded in the arXiv licence metadata for this exact version and shown on the abstract page https://arxiv.org/abs/2308.15913v4. Full licence notice: \texttt{licenses/arxiv-2308.15913-CC-BY-4.0.txt}.\par\smallskip

\noindent\textit{Editorial scope.} Counted unit: the article's cor cor (source0.tex lines 668--671). The article's complete original proof of it is delivered with it (source0.tex lines 672--677, one \texttt{proof} environment of 6 lines). No mathematics is written, repaired or re-derived: the statement and the proof are the article's own lines, in the article's own notation, and no retained line is substituted or re-flowed.  Retained uncounted as the source-local context the proof needs: lines 601--604.  Nothing was omitted from any retained span, so the package carries no omission record.  No retained line contains a citation, so the wrapper carries no bibliography and no bibliography entry.  No retained line contains a figure environment, an image-inclusion command or drawing code, so no image is delivered and the package declares no asset dependency.  Editorial formatting only, in addition: an AMS \texttt{amsart} preamble that restates the article's own declarations and macros used by the retained lines, namely the environments \texttt{thm}, \texttt{cor}, \texttt{pf} on separate counters, together with the standard packages those lines need.  The retained environments print in this document as Theorem 1, Corollary 1, Pf 1 in order of appearance, whereas the article prints the same items with the numbering of its own full document.  The complete native source of the article as distributed in the arXiv e-print for this exact version is bundled unchanged as \texttt{sources/146283/source0.tex}.
\begin{thm}\label{Rohi-Vasu-P2-Theorem-002}
Let $f:B^n\to \IC^n$ be a $(K,K')$-quasiregular mapping of the unit ball $B^n$ into $\IC^n$ with $\lambda_f(0)\geq \alpha>0$. Then
$$\beta_f\geq \frac{\alpha^2}{4(2K\alpha+K'+\alpha)}.$$ 
\end{thm}

\begin{cor}
Let $f:B^n\to \IC^n$ be a $(K,K')$-quasiregular mapping of the unit ball $B^n$ into $\IC^n$ with $|\operatorname{det} f'(0)|=\alpha>0$. Then
$$\beta_f\geq \frac{\alpha^{2/n}}{4(2K\alpha^{1/n}+K'+\alpha^{1/n})}.$$ 
\end{cor}

\begin{pf}[{\bf Proof}]
Since $$\alpha=|\operatorname{det} f'(0)|\leq \lambda_f^n(0),$$
we have,
     $$\lambda_f(0)\geq\alpha^{1/n}.$$
Now by replacing $\alpha$ by $\alpha^{1/n}$ in Theorem \ref{Rohi-Vasu-P2-Theorem-002}, we obtain the desired result.
\end{pf}

\end{document}
